\setcounter{page}{2}
\chapter*{Abstract}
\label{chapter:abstract}
Die zuverlässige Klassifikation von Anomalien in komplexen Systemen wirft eine zentrale Frage auf: Wie kann sichergestellt werden, dass ein als Anomalie identifizierter Datenpunkt tatsächlich eine Abweichung vom Normalverhalten darstellt, und nicht lediglich ein seltener, aber gültiger Ausdruck innerhalb des Systems? \\
Mit der zunehmenden Größe und Komplexität adaptiver Systeme stoßen klassische, meist lineare Verfahren zur Anomalieerkennung an ihre Grenzen. Besonders im Bereich von Big-Data-Architekturen, in denen Datenströme hochdimensional und dynamisch sind, gewinnt der Einsatz künstlicher Intelligenz zunehmend an Bedeutung.
Diese Seminararbeit beleuchtet den Einsatz generativer KI zur Anomalieerkennung in adaptiven Systemen. Anhand eines Fallbeispiels aus der aktuellen Fachliteratur wird gezeigt, wie generative adversariale Netzwerke zur Modellierung normaler Datenverteilungen genutzt werden können, um potenzielle Anomalien zu identifizieren. Ziel der Arbeit ist es, die theoretischen Grundlagen, Chancen und Herausforderungen dieser Technologie darzustellen und kritisch zu reflektieren, inwiefern generative KI einen Beitrag zur Weiterentwicklung klassischer Anomalieerkennungsverfahren leisten kann.
